\chapter{Обзор литературы}\label{ch:ch1}

Настоящая глава формирует теоретическую и методологическую базу диссертационного исследования. Её цель \--- систематизировать ключевые концепции и инструменты, необходимые для моделирования и анализа процессов фазовых превращений в рамках молекулярно-динамического подхода. 

В первом разделе даётся детальное описание метода молекулярной динамики (МД). Особое внимание уделяется физическим основам и математическому аппарату МД, а также обзору потенциалов межатомного взаимодействия, используемых в работе, поскольку их адекватный выбор является залогом достоверности получаемых результатов.

Второй раздел посвящён классической и современной проблематике фазовых переходов первого рода. Рассматривается их термодинамическая классификация, история изучения и фундаментальный механизм зарождения новой фазы — нуклеация.

Третий раздел фокусируется на поверхностных явлениях и фазовом равновесии. Подробно разбираются методы вычисления поверхностного натяжения \--- ключевого параметра, определяющего энергетику нуклеации, непосредственно из данных МД-моделирования.

Наконец, в четвёртом разделе анализируются процессы переноса. Изучается роль диффузии в кинетике роста зародышей и вводится коэффициент термической аккомодации как фундаментальная характеристика, связывающая атомистические процессы на межфазной границе с макроскопическим описанием тепло- и массообмена. 


\section{Метод молекулярной динамики}

\subsection{Общие сведения}

Метод молекулярной динамики заключается в решении уравнений Ньютона для всех $N$ частиц системы~\cite{frenkel_understanding_2001}:

\begin{align}
	m_i\frac{\overrightarrow{d^2r_i}}{dt^2} = \overrightarrow{F_i}(r_1, r_2, ... r_n),
\end{align}

где $m_i$ \--- масса частицы с номером $i$, $r_i$ \--- её координата, а $F_i$ \--- скорость. Сила рассчитывается как  градиент потенциальной энергии, взятый со знаком минус:

\begin{align}
	\overrightarrow{F_i}(r_1, r_2, ... r_n) = -\overrightarrow{\nabla}U_i(r_1, r_2, ... r_n),
\end{align}

где $U_i$ \--- потенциальная энергия $i$-ого атома. Она определяется через ряд функций, называемых потенциалом взаимодействия. Потенциал межатомного взаимодействия определяет свойства системы. 

Далее для расчета положений атомов и их скоростей в следующий шаг по времени применяется схема Верле~\cite{verlet_computer_1967}:

\begin{subequations}
    \begin{gather}
		\overrightarrow{r_i}(t+\Delta t) = \overrightarrow{r_i}(t) + \overrightarrow{v_i}(t) + \frac{\overrightarrow{F_i}(t)\Delta t^2}{2m_i},\\
		\overrightarrow{v_i}(t+\Delta t) = \overrightarrow{v_i}(t) + \frac{\overrightarrow{F_i}(t)+\overrightarrow{F_i}(t+\Delta t)}{2m_i}\Delta t.
    \end{gather}
\end{subequations}

При численном интегрировании уравнений движения в молекулярной динамике неизбежно возникают погрешности, приводящие к неконсервативности энергии системы, что подробно обсуждается в работе~\cite{norman_stochastic_2013}. Эти ошибки обусловлены отклонением численной траектории от истинной, соответствующей минимуму действия. Накопление ошибки, величина которой зависит от шага интегрирования $\Delta t$, в конечном счёте делает результаты моделирования физически некорректными. С другой стороны, уменьшение $\Delta t$ для повышения точности ведёт к пропорциональному росту вычислительных затрат. Таким образом, при выборе шага интегрирования необходимо находить компромисс между точностью и эффективностью.

Данная проблема особенно остро проявляется в системах, содержащих жёсткие степени свободы, такие как колебания химических связей или валентных углов. Период этих колебаний часто на порядки меньше характерного времени для интересующих исследователя процессов (например, диффузии или конформационных перестроек). Использование единого, достаточно малого для описания связей шага $\Delta t$ для интегрирования всей системы становится крайне неэффективным.

Для преодоления этого ограничения используется иерархический алгоритм rRESPA (reversible Reference System Propagator Algorithm)~\cite{tuckerman_reversible_1992}. Его основная идея заключается в многочастотном интегрировании разных типов взаимодействий. Алгоритм основан на разложении оператора Лиувилля на составляющие, отвечающие за быстрые (например, внутримолекулярные) и медленные (например, межмолекулярные) взаимодействия. Это позволяет использовать разные, «вложенные» шаги по времени: малый шаг $\delta t$ для быстрых сил и более крупный шаг $\Delta t = n \cdot \delta t$ (где $n$ — целое число) для медленных сил. Например, колебания связей можно интегрировать с шагом в 0.5 фс, в то время как вандерваальсовы взаимодействия — с шагом в 5 фс. Такой подход обеспечивает значительный выигрыш в вычислительной эффективности без потери точности и обратимости во времени, что является ключевым преимуществом алгоритма rRESPA.

\begin{align}
	iL = iL_1 + iL_2,
\end{align}

Для такого разложения справедливо:

\begin{align}
	e^{i(L_1+L_2)t} = (e^{i(L_1 + L_2)t/N_t})^{N_t} = (e^{i L_1 \Delta t /2} e^{i L_2 \Delta t} e^{i L_1 \Delta t /2})^{N_t} + O(t^3/P^2),
\end{align}

где $\Delta t = t/N_t$ \--- шаг по времени, а $t$ и $N_t$ \--- время моделирования и число шагов соответственно. Тогда оператор эволюции примет вид:

\begin{align}
	 e^{i L \Delta t} = e^{i L_1 \Delta t /2} e^{i L_2 \Delta t} e^{i L_1 \Delta t /2}.
\end{align}
  
Для успешного начала моделирования необходимо корректно задать начальные условия системы. На практике наиболее распространены два подхода к генерации начальных координат атомов: случайное размещение в объёме расчётной ячейки и размещение в узлах кристаллической решётки~\cite{frenkel_understanding_2001}. Каждый метод имеет свои преимущества и ограничения.

При случайном размещении существует вероятность, что атомы окажутся на критически малых расстояниях, что приведёт к возникновению огромных сил отталкивания и мгновенной нестабильности системы. Для устранения этих артефактов перед началом динамического моделирования обязательно проводится процедура энергетической минимизации (релаксации) конфигурации.

Начальная конфигурация на основе кристаллической решётки, как правило, является более физичной и лишена проблемы атомных перекрытий. Это особенно важно для систем, содержащих сложные молекулы или циклы, случайное размещение которых с последующей минимизацией может привести к их «запутыванию» и формированию артефактных конформаций. Однако главный недостаток этого подхода заключается в наложении на систему искусственной долговременной пространственной корреляции, не свойственной жидкой или газовой фазе. Для её устранения и выхода на истинное равновесное состояние может потребоваться значительно более продолжительное предварительное моделирование.

Важнейшим элементом постановки вычислительного эксперимента является выбор граничных условий. Периодические граничные условия являются стандартным выбором для моделирования объёмных свойств конденсированных сред. При их использовании расчётная ячейка мысленно повторяется бесконечное число раз во всех направлениях, создавая для атомов на её границах полное окружение, как в бесконечной среде. Это позволяет исключить поверхностные эффекты и корректно моделировать макроскопические свойства, используя ограниченное число частиц.

Однако в ряде случаев периодические граничные условия неприменимы. Например, при моделировании столкновения одиночного атома с изолированным нанокластером, как описано в Главе~\cref{ch:ch6}, использование периодических граничных условий приведёт к артефактам: один и тот же атом мог бы многократно взаимодействовать с исходным кластером или его периодическими образами, что физически некорректно и затрудняет интерпретацию данных о единичном акте соударения. Для таких систем используются фиксированные граничные условия, при которых частицы могут свободно покидать расчётную область, а взаимодействие с внешней средой отсутствует или моделируется явно.

\subsection{Потенциалы межатомного взаимодействия}

Наиболее важной характеристикой, определяющей систему, является потенциал межатомного взаимодействия. Он определяет силы, действующие между атомами в вычислительной ячейке.

Самым первым использованным потенциалом является потенциал Леннард-Джонса~\cite{Lennard-Jones_1924}:

\begin{align}
	U(\overrightarrow{r}) = 4 \epsilon (\frac{\sigma}{r}^{12}-\frac{\sigma}{r}^6),
\end{align}

где $\epsilon$ соответствует минимальной потенциальной энергии взаимодействия двух частиц, так называемой "глубине потенциальной ямы", а $\sigma$ \--- расстоянию, где силы отталкивания равны силам притяжения между двумя частицами, то есть потенциальная энергия проходит через 0.

Этот потенциал хорошо описывает свойства инертных газов, соответствует энергии взаимодействия двух наведенных диполей и очень прост в расчётах: для вычисления члена, отвечающего за отталкивание, нужно $(\sigma / r)^6$ возвести в квадрат. Также этот потенциал используется для расчета межмолекулярных взаимодействий в углеводородах. Например, в моделях TraPPE-UA, OPLS и COMPASS. 

Однако, не для всех систем подходит такое описание, и разрабатываются иные потенциалы межатомного взаимодействия. Например, в потенциале Букингема~\cite{buckingham_1938} изменен член, отвечающий за отталкивание, на экспоненциальный:

\begin{align}
	U(\overrightarrow{r}) = A e^{-B r} - \frac{C}{r}^6),
\end{align}

где $A, B, C$ \--- константы. 

Наряду с потенциалом Букингема в таких задачах, как взаимодействие атомов металлов с инертным газом, используется потенциал Морзе~\cite{morse_diatomic_1929}. В нем оба слагаемых заменены на экспоненциальные функции:

\begin{align}
	U(\overrightarrow{r}) = D [e^{-2 \alpha (r - r_0)} - 2 e^{-\alpha (r - r_0)}],
\end{align}

где $D$ \--- глубина потенциальной ямы, $r_0$ \--- расстояние между атомами, на котором находится минимум потенциальной энергии, $\alpha$ \--- коэффициент.

Все эти потенциалы межатомного взаимодействия учитывают только расстояния между парами атомов. Такой способ расчета потенциальной энергии плохо подходит, например, для металлов, для которых необходимо учитывать общую электронную плотность. Для этого существуют многочастичные потенциалы, наподобие EAM(embedded atom method)~\cite{daw_embedded-atom_1984}, в которых взаимодействие между атомами зависит также от величины электронной плотности, в которую вносят вклад все соседние атомы. Эта величина, как и все остальные, необходимые для данного потенциала, задаются таблично. Таким образом, формула для EAM принимает вид:

\begin{align}
	U(\overrightarrow{r}) = U_{pair}(r) + U_{many}(\rho),
\end{align}

где $U_pair$ \--- энергия парного взаимодействия, $U_{many}$ \--- энергия, зависящая от электронной плотности.

Данный тип потенциалов получил широкое развитие, поскольку достаточно хорошо описывают поведение металлов. В частности, можно отметить потенциал Финниса-Синклера~\cite{ackland_computer_1997}, который хорошо описывает плавление металлов и отличается тем, что учитывает тип атомов, которыь е вносят вклад в электронную плотность:

\begin{align}
	U(\overrightarrow{r}) = U_{pair}(r) + U_{\alpha \beta}(\rho),
\end{align}

где $U_{\alpha \beta}$ \--- энергия многочастичного взаимодействия, зависящая от типов атомов.


Однако, взаимодействия между атомами могут быть не только за счет сил Ван-дер-Ваальса. Внутримолекулярные взаимодействия описываются сходным образом. Существует множество моделей, описывающих молекулы. Чаще всего в них энергия взаимодействия  атомов разбивается на энергии связей, углов и двугранных углов:

\begin{align}
	U(r) = U_{bond} + U_{angle} + U_{dihedral}
\end{align}

Существует множество способов описания этих взаимодействий. Например, для связей часто применяется гармонический потенциал:

\begin{align}
	U(r) = k (r - r_0)^2,
\end{align}

где $r$ \--- длина связи, $r_0$ \--- равновесная длина связи, $k$ \--- коэффициент жесткости. Следует отметить, что коэффициент жесткости, например, в углеводородах достаточно велик, что вынуждает рассчитывать связи по схеме RESPA с меньшим шагом по времени.

Аналогично рассчитываются энергии углов. В случае с двугранными углами используются, например, суммы Фурье:

\begin{align}
	U(\phi) = \sum_{i=1, m}[K_i(1+cos(n_i\phi - d_i))],
\end{align}

В случае модели COMPASS имеется вклад энергии связей в энергию углов и двугранных углов, что будет более подробно показано в главе~\cref{ch:ch2}.

В некоторых случаях, например, в модели воды TIP4P~\cite{abascal_general_2005}, применяются жесткие связи и углы. Это означает, что длины связей и величины углов остаются неизменными. Это достигается, например, с помощью алгоритма SHAKE~\cite{ryckaert_numerical_1977}. 

\subsection{Дальнодействие и производительность}

Очевидно, что выбор потенциала межатомного взаимодействия влияет на скорость проведения расчетов. Одним из основных факторов, влияющих на производительность, является радиус обрезания. Это расстояние, начиная с которого энергия взаимодействия между атомами приравнивается к нулю. Например, для потенциала Леннард-Джонса:

\[
U(r) = \left\{
\begin{alignedat}{2}
    &4 \epsilon (\frac{\sigma}{r}^{12}-\frac{\sigma}{r}^6), \quad &\text{eсли } r\geqslant r_{cut} \\
    &0, \quad & \text{eсли } r_{cut}<0
\end{alignedat}
\right.
\]

Для радикального сокращения вычислительных затрат с  $O(N^2)$  до  $O(N \log N)$ , где  $N$  — число частиц, широко применяется метод списков соседей. Его суть заключается в предварительном составлении для каждого атома перечня всех потенциальных партнёров по взаимодействию, находящихся в пределах заранее заданного радиуса  $r_{\text{list}}$ . Критически важно, что этот радиус, как правило, превышает радиус обрезания потенциала  $r_{cut}$ . Разность  $r_{\text{list}} - r_{cut}$  формирует так называемый скин-слой. Атомы, попавшие в этот слой, включаются в список, но непосредственно не взаимодействуют на текущем шаге. Такой подход позволяет реже (например, раз в 10-20 шагов) выполнять трудоёмкую процедуру перестроения списков, исключая при этом артефакты, когда частица за один шаг могла бы войти в зону взаимодействия  $r_{cut}$, не будучи учтённой. Это даёт существенный выигрыш в общей производительности симуляции.

Однако обрезание потенциала на конечном радиусе $r_{cut}$ приводит к систематическому занижению энергии и других термодинамических величин системы, поскольку игнорируется вклад от части взаимодействий на расстояниях  $r > r_{cut}$. Для коррекции этих ошибок существует несколько стандартных подходов, подробно рассмотренных в литературе (например, в ~\cite{sun_compass_1998}).

Для решения этой проблемы применяется несколько методов. Один из них \--- метод сдвига потенциала, который заключается в добавке к потенциалу константы, делающей функцию непрерывной в $r = r_{cut}$~\cite{toxvaerd_communication_2011}:

\[
U(r) = \left\{
\begin{alignedat}{2}
    &U_{lj}(r) - U_{lj}(r_{cut}), \quad &\text{eсли } r\geqslant r_{cut} \\
    &0, \quad & \text{eсли } r_{cut}<0, 
\end{alignedat}
\right.
\]

где $U_{lj}(r) = 4 \epsilon (\frac{\sigma}{r}^{12}-\frac{\sigma}{r}^6)$. Эта процедура устраняет разрыв в значении самой функции, однако производная потенциала (а следовательно, и сила) в точке обрезания  $r_{cut}$  остаётся разрывной. Это приводит к нарушению сохранения энергии и может вызывать нефизичные артефакты в динамике системы, особенно при частых пересечениях атомами границы $r_{cut}$.

Для обеспечения гладкости как потенциала, так и силы в точке обрезания применяется метод "непрерывной сшивки". Его идея заключается в замене исходного потенциала на сглаживающую функцию (как правило, многочлен) в промежуточной области между двумя радиусами:  $r_{\text{on}}$  и  $r_{\text{off}}$  ( $r_{\text{on}} < r_{\text{off}}$)~\cite{marrink_coarse_2004}.

Функция сшивки $S(r)$ конструируется таким образом, чтобы удовлетворять четырём граничным условиям:

\begin{enumerate}

\item При  $r = r_{\text{on}}$:  $S(r_{\text{on}}) = U_lj(r_{\text{on}})$  и  $S'(r_{\text{on}}) = U'(r_{\text{on}})$, что обеспечивает плавное сопряжение с исходным потенциалом.
\item При  $r = r_{\text{off}}$:  $S(r_{\text{off}}) = 0$  и  $S'(r_{\text{off}}) = 0$, что обеспечивает плавный переход к нулевому значению.

\end{enumerate}

Такой подход гарантирует, что первая производная потенциала (сила) является непрерывной во всём диапазоне расстояний, что критически важно для корректного численного интегрирования уравнений движения и сохранения полной энергии системы. 

Для потенциала Леннард-Джонса и многих других потенциалов межатомного взаимодействия свойственна сходимость следующего интеграла вне сферы конечного радиуса:

\begin{align}
	\label{eq:integral}
	U_{tail} = \int_{r_{cut}}^{\inf} \rho U(r) 4 \pi r^2 dr,
\end{align}

где $\rho$ \--- плотность, $U(r)$ \--- рассматриваемый потенциал межатомного взаимодействия.

При помощи подобного интеграла рассчитываются дальнодействующие поправки к потенциалу межатомного взаимодействия. Это позволяет уменьшить ошибку, возникающую из-за наличия радиуса обрезания. 

У этого метода есть ряд недостатков. Например, он не подходит для гетерогенной среды. Если моделируется две или более фазы, плотности в них будут разными, что противоречит предположению, сделанному в формуле.

Для взаимодействий по закону Кулона такой метод также не является подходящим, поскольку интеграл выше~\cref{eq:integral} на бесконечности сходится условно. Потому используются, например, потенциалы для плазмы, такие как потенциал Юкавы~\cite{yukawa_interaction_1955}:

\begin{align}
    U(r) = - \frac{Z_1 Z_2 q_e^2}{k_1} e^{-k_2r},
\end{align}

где $k_1, k_2$ \--- некоторые константы, $Z_1, Z_2$ \--- зарядовые числа взаимодействующих частиц, $q_e$ \--- заряд электрона. Этот потенциал описывает экранирование, которое возникает в плазме, за счет обратного экспоненциального множителя. Такой потенциал является интегрируемым на бесконечности, потому для него возможно рассчитать поправки, компенсирующие отсутствие дальнодействующей части.

Ещё одним способом учесть дальнодействие в системе является метод Эвальда~\cite{kolafa_cutoff_1992}, который иногда, например, в GROMACS, реализован и для потенциалов, обычно рассчитываемых через поправки. Этот метод заключается в помещении заряда в оболочку противоположного знака. Далее происходит разложение в ряд Фурье для заряда в оболочке, которая задается формулой:

\begin{align}
    \rho_{Gauss}(r) = - q_i (\alpha / \pi)^{3/2} e^{-\alpha r^2},
\end{align}


где $\rho_{Gauss}$ \--- электронная плотность, задаваемая гауссианой, $q_i$ \--- экранируемый заряд, $\alpha$ \--- коэффициент, определяемый из соображений вычислительной эффективности.

В таком случае дальнодействующая часть будет равна:

\begin{align}
    U_1 = \frac{1}{2}\sum_{i}q_i \phi(r_i) = \frac{1}{2V}\sum_{k \neq 0} \frac{4\pi}{k^2}\rho(\overrightarrow{r})e^{-k^2/4\alpha},
\end{align}

где $\phi(r_i)$ \--- значение потенциала кулоновского потенциала в точке, где располагается $i$-ый атом, $k = 2\pi/L$, где $L$ длина кубической вычислительной ячейки.

Однако, это выражение включает взаимодействие точечного заряда и окружающей его гауссианы. Потому требуется вычесть следующее выражение:

\begin{align}
    U_{self} = (\alpha/\pi)^{1/2} \sum_{i = 1}^N q_i^2.
\end{align}

Таким образом, выражение для полной энергии кулоновского взамиодействия, включающее в себя и короткодействующую часть, примет вид:

\begin{equation}
    U_{self} = \frac{1}{2}\sum_{i}q_i \phi(r_i) = \frac{1}{2V}\sum_{k \neq 0} \frac{4\pi}{k^2}\rho(\overrightarrow{r})e^{-k^2/4\alpha} - (\alpha/\pi)^{1/2} \sum_{i = 1}^N q_i^2 + \frac{1}{2}\sum_{i \neq j}^N \frac{q_i q_j erfc(\alpha r_{ij})}{r_{ij}}.
\end{equation}

Следует заметить, что это требующая значительных ресурсов по сравнению с обычным расчетом дальнодействующих поправок процедура, потому был найден способ её ускорить. Для этого создается решетка, на которую проецируются имеющиеся в системе заряды. Этот метод называется particle-particle particle-mesh(PPPM)~\cite{eastwood_p3m3dp-three-dimensional_1984}.

Ещё один метод учета дальнодействия в системе показан в работе~\cite{Смирнов__2022}. Этот метод основан на моделировании системы с несколькими радиусами обрезания и экстраполяции полученных результатов на бесконечность. Такой метод показал хорошее согласие с экспериментом и достаточно высокую по сравнению с расчетом методом Эвальда производительность.

\subsection{Термостаты и баростаты}

Моделирование методом молекулярной динамики позволяет создавать системы, фиксируя различные условия. Самым простым вариантом является моделирование закрытого сосуда в теплоизолирующей оболочке, NVE-ансамбль. В этом случае фиксированны количество частиц, объем и энергия. 

Существует возможность зафиксировать температуру при моделировании, получив таким образом NVT-ансамбль, в котором происходит обмен энергией с окружающей средой, но не массой. Для этого потребуется усложнить алгоритм, добавив в систему термостат. Существует три самых распространенных вида термостатов: Берендсена~\cite{berendsen_molecular_1984}, Ланжевена~\cite{allen_computer_2017} и Нозе-Гувера~\cite{nose_unified_1984}.

Термостат Берендсена~\cite{berendsen_molecular_1984} является простейшим из перечисленных:

\begin{align}
    \frac{dT}{dt} = \frac{T-T_0}{\tau},
\end{align}

где $T$ и $T_0$ \--- температура системы и заданная температура соответсвенно, $t$ \--- время, $\tau$ \--- параметр, определяющий силу термостата. Его влияние заключается в следующем: каждая молекула системы, в зависимости от её текущей скорости, подвергается корректировке с использованием заданного коэффициента термализации. Это позволяет "отбирать" или "добавлять" энергию в систему, чтобы поддерживать температуру на заранее определённом уровне. 

Одним из недостатков этого термостата является нарушение динамики за счет изменения скоростей не вполне естественным образом. Также в результате работы этого термостата не устанавливается корректное распределение Гиббса.

Термостат Ланжевена~\cite{allen_computer_2017} также иногда называется стохастическим за счет добавления в него случайной силы:

\begin{align}
    m\frac{dv}{dt} = \frac{dU}{dr} - \lambda v + \xi(t),
\end{align}

где $v$ \--- скорость, $\lambda$ \--- коэффициент трения, $\xi(t)$ \--- некоторая функция для создания шума в системе. Коэффициент трения приводит к экспоненциальному уменьшению температуры, а шум подбирается так, чтобы приближаться к нужной температуре. Стохастичность уравнений движения, однако, принуждает переходить в систему отсчета, связанную с центром масс системы, чтобы избежать случая, когда частицы перестают двигаться относительно друг друга, но приобретают значительные сонаправленные скорости. Также этот термостат влияет на значения коэффициентов переноса.

Одним из наиболее широко используемых термостатов является термостат Нозе-Гувера~\cite{nose_unified_1984}, основанный на внесении поправок в лагранжиан системы. Новый лагранжиан принимает вид:

\begin{align}
    L_{Nose} = \sum_{i=1}^{N} \frac{1}{2} m_i s^2 \dot{\vec{r_i^2}} - U(\vec{r}^N) _ \frac{1}{2}Q \dot{s}^2 - \frac{3N}{kT} ln s,
\end{align}

где $s$ \--- параметр, характеризующий силу термостата, $Q$ \--- масса, которая к нему относится. 

В таком случае выражения для импульсов меняют вид:

\begin{align}
    \vec{p_i} = \frac{\partial L}{\partial \dot{\vec{r_i}}} = m_i s^2 \dot{\vec{r_i}}
\end{align}

Для новой координаты также определяется импульс:

\begin{align}
    p_s = \frac{\partial L}{\partial \dot{s}} = Q \dot{s}
\end{align}

Тогда гамильтониан системы примет вид:

\begin{equation}
    H_{Nose} = \sum_{i=1}^N \frac{\vec{p_i^2}}{2m_is^2} + U(\vec{r}^N) + \frac{p_s^2}{2Q} + \frac{3N ln s}{kT}
\end{equation}

Можно представить новые уравнения движения как уравнения движения с переменным шагом по времени. $s$ служит коэффициентом масштабирования между реальным шагом и тем, с которым происходит интегрирование. Сами уравнения движения примут вид:

\begin{subequations}
	\begin{gather}
		\frac{d\vec{r_i}}{dt} = \frac{\partial H_{Nose}}{\partial \vec{p_i}} = \frac{\vec{p_i}}{m_i s^2} \\
		\frac{d\vec{p_i}}{dt} = -\frac{\partial H_{Nose}}{\partial \vec{r_i}} = \frac{\partial U(\vec{r}^N)}{\partial \vec{r_i}} \\
		\frac{ds}{dt} = \frac{\partial H_{Nose}}{\partial p_s} = \frac{p_s}{Q} \\
		\frac{d p_s}{dt} = -\frac{\partial H_{Nose}}{\partial s} = \frac{1}{s} \sum_{i=1}^N (\frac{\vec{p_i}^2}{m_i s^2} - \frac{3N}{kT}) 
	\end{gather}
\end{subequations}

Недостатком термостата Нозе-Гувера являются нефизичные флуктуации температуры. Для уменьшения их влияния используется цепочка из таких термостатов, применяемых последовательно.

Ещё одной системой, поддающейся моделированию, является вещество, находящееся под поршнем внутри сосуда с проводящими тепло стенками. В такой системе регулируемыми параметрами оказываются число частиц, давление и температура. Это NPT-ансамбль. Давление в молекулярной динамике определяется следующим образом:

\begin{equation}
    P = \frac{1}{dV}[\sum_{i=1}{N}(\frac{p_i^2}{m_i} + \vec{r_i} \vec{F_i}) - dV]
\end{equation}

Для реализации такого случая требуется баростат. Одним из наиболее распространенных баростатов является баростат Шиноды, который состоит из цепочки Нозе-Гувера~\cite{Shinoda}. Изменение давления происходит за счет объема вычислительной ячейки. Такой баростат эффективен для систем с резкими изменениями плотности как, например, в системах с фазовыми переходами. Он также сохраняет фазовый объем.

Существует также возможность моделирования открытой системы, однако в рамках данной диссертации в таких моделях не возникло необходимости, потому они описаны не будут.

Молекулярное моделирования для решения поставленных задач было реализовано на оснвое открытого пакета молекулярной динамики Large-scale Atomic/Molecular Massively Parallel Simulator(LAMMPS)~\cite{thompson_lammps_2022}.


\section{Фазовые переходы 1-го рода и влияющие на них факторы – массообмен, теплообмен и поверхностное натяжение}

Фазовые переходы, то есть скачкообразные изменения свойств вещества при плавном изменении внешних параметров (температуры, давления), представляют собой фундаментальное явление в физике конденсированного состояния. С термодинамической точки зрения, фазовые переходы традиционно классифицируются по работам Ландау~\cite{landau5} на переходы первого и второго рода. Критерием классификации является поведение термодинамических потенциалов, таких как свободная энергия Гиббса $G$ .

При фазовом переходе первого рода скачком изменяются первые производные термодинамического потенциала, которые соответствуют таким экстенсивным параметрам, как энтропия  $S = -\left(\frac{\partial G}{\partial T}\right)_P$  и объём  $V = \left(\frac{\partial G}{\partial P}\right)_T$. Это означает, что переход сопровождается выделением или поглощением скрытой теплоты (теплота плавления, испарения) и скачкообразным изменением плотности. Типичными примерами служат плавление, кипение и сублимация.

В отличие от этого, при фазовом переходе второго рода первые производные термодинамического потенциала изменяются непрерывно (нет скачка плотности и выделения скрытой теплоты), однако разрыв испытывают вторые производные, такие как теплоёмкость  $C_P$, сжимаемость и термический коэффициент расширения. Примером является переход в сверхпроводящее состояние.

Ключевой особенностью фазовых переходов первого рода является их склонность к затруднённому протеканию в метастабильных условиях. Поскольку сосуществующие фазы при таком переходе имеют различную плотность и структуру, возникновение новой фазы внутри исходной не может произойти одновременно во всём объёме. Вместо этого система проходит через образование локальных неустойчивостей — мелких зародышей или кластеров новой фазы. Этот начальный стадийный процесс, известный как нуклеация (зародышеобразование), и является кинетическим «бутылочным горлышком» перехода.


Процесс перехода между фазами можно описать следующей последовательностью:

\begin{enumerate}

\item Нуклеация: в переохлажденной или пересыщенной исходной фазе вследствие флуктуаций спонтанно образуются наноразмерные зародыши новой фазы.
\item Рост: те зародыши, размер которых превысил некий критический порог, становятся устойчивыми и начинают расти, потребляя вещество из исходной фазы.
\item Коалесценция: растущие домены новой фазы сливаются друг с другом, формируя общую структуру.
\item Завершение перехода: новая фаза полностью замещает исходную.

\end{enumerate}


Наиболее устоявшимся и фундаментальным подходом к описанию стадии нуклеации является классическая теория нуклеации (КТН), разработанная в работах Фольмера, Вебера, Беккера и Дёринга~\cite{pruppacher_microphysics_1996}. В её основе лежит концепция критического зародыша — кластера новой фазы, находящегося в неустойчивом равновесии с исходной фазой. Его существование является следствием конкуренции двух факторов: объёмной свободной энергии, выигрышной при образовании более стабильной фазы, и поверхностной свободной энергии, связанной с созданием границы раздела фаз и являющейся энергетическим барьером для процесса.

Таким образом, КТН предоставляет количественную модель для расчёта скорости образования зародышей, что является центральной задачей при описании кинетики фазовых переходов первого рода.

\subsection{Классическая теория нуклеации}

Особенностью фазовых переходов первого рода является существование метастабильных состояний. Это означает, что система не переходит в новую фазу мгновенно при достижении линии фазового равновесия. Она может достаточно долго существовать в исходном состоянии, такаом, как переохлажденный пар, пока не образуется достаточная неоднородность условий для фазового перехода. 

Такой неоднородностью служат стабильные зародыши или кластеры новой фазы. Процесс их образования называется нуклеацией или зародышеобразованием. Этот процесс служит кинетическим бутылочным горлышком для фазового перехода, поскольку её скорость экспоненциально зависит от высоты энергетического барьера, необходимого для образования новой фазы.

Рассмотрим природу этого барьера, для чего нужно проанализировать термодинамический потенциал Гиббса, необходимый для образования зародыша новой фазы с радиусом $r$ внутри исходной фазы. Изменение этого потенциала состоит из двух слагаемых. Первое из них показывает уменьшение энергии с ростом объёма новой фазы, поскольку в условиях пересыщения для вещества выгоднее находиться в новой фазе. Второе выражает энергию формирования поверхности и зависит от поверхностного натяжения:

\begin{equation}

\Delta G = \frac{4}{3} \pi r^3 \Delta g_v + 4 \pi r^2 \sigma,

\end{equation}

где $\Delta G$ \--- изменение термодинамического потенциала Гиббса при формировании новой фазы, $\Delta g_v$ \--- плотность изменения потенциала Гиббса без учета влияния поверхности, $\sigma$ \--- поверхностное натяжение. 
 
У этой функции есть максимум при некотором значении $r$:

\begin{equation}

r* = 2\frac{\sigma}{\Delta g_v},

\end{equation}

Это значение $r$ называется критическим радиусом $r_{crit}$, а зародыш такого размера \--- критическим зародышем. Можно увидеть, что для размеров меньше критического работа по увеличению кластера положительна, а в противном случае \--- отрицательна. Таким образом, малые кластеры склонны распадаться. В случае возникновения достаточно большого зародыша новой фазы рост становится более энергетически выгодным, чем распад, и вероятен сценарий дальнейшего увеличения объема новой фазы.  

Следовательно, чем меньше размер критического кластера, тем выше вероятность образования стабильного зародыша новой фазы. На эту величину влияет множество факторов, среди которых степень пересыщения, состав системы и поверхностное натяжение.  


\section{Поверхностные явления и их описание}

Один из способов понизить высоту барьера потенциальной энергии \--- изменить состав системы. Гетерогенная нуклеация протекает при более низких степенях пересыщения~\cite{prutton_surface_2000}. Первый подраздел рассматривает вопрос адсорбции на границе межфазной поверхности в двухкомпонентной системе.

Наличие в системе других компонентов снижает поверхностное натяжение. Например, в работе~\cite{bhatt_molecular_2004} исследуется влияение солей на поверхностное натяжение воды с помощью метода молекулярной динамики. Во втором подразделе дан обзор работ, посвященных поверхностному натяжению в многокомпонентных системах.

\subsection{Поверхность Гиббса. Адсорбция}

В данной работе рассмотрены фазовые переходы для двух различных систем: слоя жидкости, окруженного паром, и уединенного кластера, на который налетают атомы. Поверхностные явления исследуются для первого случая с использованием теории Гиббса~\cite{kalikmanov_nucleation_2013}. Поверхностью в рамках данной теории является плоская поверхность с нулевой толщиной. Потому считается, что ячейка разделена на объёмы, соответствующие жидкой и газовой фазам:

\begin{equation}
    V = V_{liq} + V_{gas}
\end{equation}

Рассматриваются равновесные случаи, то есть выполняются все три условия равновесия: равенство давлений, температур и химических потенциалов. Для числа частиц при этом требуется учесть поверхностный избыток, который и показывает величину адсорбции:

\begin{equation}
    N = N_{liq} + N_{gas} + N_{exc}
\end{equation}

Здесь $N_{exc}$ \--- избыточное количество вещества. Поверхностная плотность избыточного количества вещества называется адсорбцией:

\begin{equation}
    \Gamma = N_{exc} / S,
\end{equation}

где $S$ \--- площадь поверхности.  

Выбор положения поверхности, очевидно, изменяет значение адсорбции для каждого из компонентов. Однако, в секции с результатами будет показано, что качественно картина не меняется. То есть изотерма адсорбции описывается уравнением того же вида независимо от способа задания поверхности и её положения. Одним из наиболее удобных способов задания поверхности является эквимолярная поверхность \--- поверхность, для которой $N_{exc} = 0$. Однако, в случае наличия нескольких компонентов это равенство не может выполняться для всех компонентов и для смеси одновременно. Следовательно, в случае смеси приходится выбирать либо компонент, для которого адсорбция будет равна нулю, либо выбирать поверхность так, чтобы суммарная адсорбция была нулевой.

Однако, плоская поверхность \--- предполагаемый случай, а реальный профиль плотности выглядит иначе. Наиболее распространенные способы аппроксимации \--- гиперболический тангенс и функция ошибок:

\begin{align}
\rho(x) = \rho_g + \frac{\rho_l-\rho_g}{2}(tanh(\frac{x-x_0}{D})+1),
\end{align}

\begin{align}
\rho(x) = \rho_g + \frac{\rho_l-\rho_g}{2}(erf(\frac{x-x_0}{D})+1),
\end{align}

где $rho_g$ и $rho_l$ \--- плотности пара и жидкости соответсвенно, $x_0$ \--- положение эквимолярной поверхности, $D$ \--- харатерная толщина поверхности.


Ключевым преимуществом прямого атомистического моделирования является возможность прямого расчёта величины адсорбции. Это создаёт основу для прямой верификации теоретических моделей, описывающих адсорбцию на межфазных границах, и позволяет оценивать адекватность приближений, закладываемых в эти модели.
Одним из важных достоинств молекулярного моделирования является возможность прямого расчёта величины адсорбции. Это создаёт основу для прямой верификации теоретических моделей, описывающих адсорбцию на межфазных границах, и позволяет оценивать адекватность приближений, закладываемых в эти модели. Более того, оно позволяет визуально наблюдать адсорбцию на зависимости плотности от координаты, ось которой направлена перпендикулярно рассматриваемой поверхности. Она выглядит как пик по сравнению с плотностями в обеих фазах. 

Адсорбция часто рассматривается экспериментально. В качестве примеров можно привести работу Байдакова и др.~\cite{baidakov_surface_2013}. В ней изучается смесь метан-этан и измеряются экспериментально величины плотностей фаз и поверхностного натяжения на межфазной границе. Далее на основе этих данных в предположении идеальности раствора, а также что пар над ним \--- идеальный газ. В работе~\cite{baidakov_Ar_2006} экспериментально исследуются растворы гелий-аргон и неон-аргон. В ней были получены зависимости поверхностного натяжения смесей от давления и на основе теории идеальных растворов оценена величина адсорбции гелия в поверхностном слое в зависимости от его доли в жидкой фазе. 


\subsection{Поверхностное натяжение}

Одной из других интересных характеристик поверхности является поверхностное натяжение. Оно показывает велчичну свободной энергии, необходимую для увеличения площади поверхности на одну единицу:

\begin{equation}
    \sigma = (\frac{\partial F}{\partial S})_{N, V, T} = (\frac{\partial G}{\partial S})_{N, P, T},
\end{equation}

Как было замечено выше, величина поверхностного натяжения напрямую влияет на высоту потенциального барьера образования кластера новой фазы. Таким образом, для систем, для которых важен фазовый состав, например, для газоконденсата в месторождении, где он находится в пористой породе. Увеличение поверхностного натяжения значительно влияет на давление в порах и может помешать извлечь полезные ископаемые.  

В литературе тема поверхностного натяжения освещена достатоно широко. В обзоре приводит анализ определений поверхностного натяжения и теоретико-вычислительных методов  

В научной литературе проблеме определения поверхностного натяжения и разработке соответствующих вычислительных методик уделяется значительное внимание. Так, в обзоре~\cite{Tovbin_2018} представлен систематический анализ существующих определений данной величины и сопоставлен широкий спектр теоретических и вычислительных подходов к её расчёту. Развитие градиентной теории, в частности, позволило получить для смеси метана и азота не только значения поверхностного натяжения, но и величины адсорбции, а также параметры профилей плотности в межфазной области~\cite{baidakov_surface_2016}. В той же работе для описания экспериментальных данных была успешно применена модель идеального раствора. Экспериментально поверхностное натяжение определяется, например, по кривизне поверхности в капилляре~\cite{kaverin2006}. 

Температурное поведение поверхностного натяжения, в том числе его зависимость от кривизны поверхности раздела, исследуется с помощью различных моделей. Например, в работе~\cite{tovbin_2022} в рамках модели решеточного газа проанализирована температурная зависимость этой величины для плоской поверхности и её изменение с размером капли при постоянной температуре. Наряду с поверхностными свойствами, методами компьютерного моделирования определяются и границы устойчивости фаз; иллюстрацией этому служит реализация алгоритма VT-flash на языке программирования Julia, представленная в исследовании~\cite{zakharov_cubiceosjl_2023}. Существуют также примеры применения атомистического моделирования для определения поверхностного натяжения жидкостей, находящихся в нанопорах~\cite{Rivera_2015}.


\section{Процессы переноса}

Помимо термодинамического барьера, определяемого поверхностным натяжением, на скорость фазового перехода существенное влияние оказывают кинетические процессы, такие как диффузия и теплоперенос.

Диффузия отвечает за подачу частиц из объёма метастабильной фазы к поверхности растущего зародыша. Однако собственно рост кластера определяется не только скоростью их поступления, но и эффективностью отвода избыточного тепла, выделяющегося при фазовом превращении. При присоединении к зародышу частица переходит в состояние с более низкой свободной энергией, и разница энергий выделяется в виде тепла. В результате растущий кластер оказывается локально перегретым относительно окружающей среды. Это повышение температуры, в соответствии с принципом Ле Шателье, создаёт термодинамическую движущую силу, стремящуюся растворить кластер, и, как следствие, замедляет его дальнейший рост.

Таким образом, возникает конкуренция двух процессов: диффузионного подвода частиц, стремящегося ускорить рост, и теплового барьера, его ограничивающего. Ключевую роль в разрешении этого противоречия играет процесс тепловой аккомодации — эффективной передачи избыточной энергии от перегретого зародыша частицам окружающей среды. Эффективность этого процесса количественно описывается коэффициентом термической аккомодации.

В данной работе процесс термической аккомодации исследован для малых кластеров железа, для которых классические макроскопические подходы могут оказаться неприменимы.

\subsection{Самодиффузия углеводородов}

После образования стабильного кластера новой фазы его рост зависит в значительной степени от диффузии. В данной работе диффузия рассмотрена для двух углеводородов, 2,2,4-триметилпентан$(C_8H_{18})$ и 1,1-дифенилэтан$(C_{14}H_{14})$~\cite{kim_correlation_2020}, рассчитанная с помощью метода молекулярной динамики.

Метод молекулярной динамики является мощным инструментом, дополняющим экспериментальные исследования, в частности, для прогнозирования транспортных свойств веществ в условиях, труднодостижимых на практике~\cite{brazhkin2018, ewen_advances_2018, maginn_best_2018}. Точность МД-расчетов критически зависит от адекватности выбранных потенциалов межатомного взаимодействия и методов последующего анализа данных.

Для вычисления транспортных свойств в методе молекулярной динамики используются равновесные и неравновесные подходы. К равновесным относятся метод Эйнштейна-Смолуховского (и его обобщение Хелфанда~\cite{helfand_transport_1960}) для расчета коэффициента диффузии и метод Грина-Кубо~\cite{green_markoff_1954, kubo_statistical-mechanical_1957}, который связывает транспортные коэффициенты с интегралами от функций автокорреляции соответствующих потоков. Оба метода, несмотря на свою строгость, могут сталкиваться с проблемами сходимости, для решения которых разрабатываются специальные вычислительные методики~\cite{kondratyuk_self-consistent_2016, zhang_reliable_2015, nichols_fourier_2015}.
 
Предсказательная сила классических МД-силовых полей, особенно при высоких давлениях, является ключевым фактором. Проведенные исследования демонстрируют возможности различных потенциалов (OPLS-AA, united-atom и др.) в воспроизведении экспериментальных данных по диффузии для н-алканов и их смесей в широком диапазоне давлений — вплоть до 1–2 ГПа~\cite{jadhao_2017, jadhao_2019, maginn_best_2018, feng_md_2013, prentice_experimental_2020, ewen_comparison_2016, glova_toward_2019, santos_self-diffusion_2019}.

Важной площадкой для объективной проверки точности МД-методов стали «Соревнования по моделированию свойств промышленных жидкостей» (Industrial Fluid Properties Simulation Challenges). В рамках 10-го и 11-го этапов ставилась задача «слепого» прогнозирования транспортных свойств различных углеводородов при давлениях до 1000 МПа~\cite{industrial_nodate}. Успешное участие в таких соревнованиях требует предварительной тщательной валидации методик на веществах с известными экспериментальными данными, такими как 2,2,4-триметилпентан (C₈H₁₈) и 2,2,4-триметилгексан (C₉H₂₀), для которых существует детальная информация по зависимости вязкости от давления вплоть до 1 ГПа~\cite{dymond_transport_1980, toraman_impact_2023}.

Таким образом, современный метод молекулярной динамики предоставляет хорошо зарекомендовавшие себя подходы для расчета транспортных свойств, а его точность постоянно проверяется и верифицируется в ходе как прямых сравнений с экспериментом, так и в рамках независимых слепых тестов.

\subsection{Термическая аккомодация}

Термическая аккомодация \--- процесс выравнивания температуры кластера с окружающей средой. При высоких степенях пересыщения она уравновешивает перегрев, и служит бутылочным горлышком процесса~\cite{eremin_influence_2008}. Этот процесс удобно рассматривать экспериментально на наночастицах металлов, поскольку для них при нормальных условиях достижимы высокие степени пересыщения, а они сами являются перспективным объектом для использования в промышленности.

Наночастицы металлов представляют значительный интерес для современных технологий благодаря их уникальным каталитическим, оптическим и механическим свойствам~\cite{gurentsov_analysis_2013, vorontsov_multiscale_2013, golovin_nanosuspension_2013}. Их применение варьируется от катализаторов синтеза углеродных наноматериалов~\cite{gurentsov_analysis_2013} до модификаторов функциональных материалов и систем адресной доставки лекарств~\cite{vorontsov_multiscale_2013, golovin_nanosuspension_2013}. Ключевым аспектом, определяющим эти свойства, является размер частиц, что делает управление процессом их синтеза первостепенной задачей.

Одним из распространенных методов получения наночастиц является конденсация из металлического пара, приводящая к образованию ансамбля кластеров с некоторым распределением по размерам. Поскольку свойства конечного продукта критически зависят от этого распределения, необходимы точные методы in-situ диагностики.

Для определения размеров наночастий в процессе их синтеза используются различные методы, такие как просвечивающая электронная микроскопия (TEM) и метод лазерно-индуцированного инкандесценции с временным разрешением (Ti-Re LII)~\cite{gurentsov_analysis_2013, eremin_nanoparticle_2006}. Последний основан на измерении кинетики остывания наноразмерных частиц после их нагрева коротким лазерным импульсом: скорость остывания, регистрируемая по спаду теплового излучения, зависит от размера частиц~\cite{gurentsov_analysis_2013}.

Данный метод успешно применяется для анализа широкого класса наноматериалов, включая металлические и углеродные частицы~\cite{gurentsov_analysis_2013, eremin_nanoparticle_2006, baranyshyn_heat_2010, gurentsov_nanoparticles_2017, nemchinsky_simulation_2024} и металлические пленки~\cite{luo_molecular_2025}.  Однако его ключевым ограничением является зависимость точности от знания эффективности теплообмена между частицей и окружающим газом, что описывается коэффициентом термической аккомодации (КТА). На практике КТА часто определяют калибровкой, сравнивая размеры, полученные методами Ti-Re LII и TEM~\cite{eremin_influence_2008}, что подчеркивает потребность в независимых способах его определения.

Коэффициенты термической аккомодации необходимы и для моделей течения разреженных газов вблизи твердых поверхностей~\cite{yousefi-nasab_review_2024}. Они определяют скорость скольжения, напряжение сдвига и температурный скачок на границе. 

Прямое компьютерное моделирование представляет собой альтернативный путь исследования коэффициента термической аккомодации и родственных процессов на атомарном уровне. В ряде работ методом молекулярной динамики (МД) изучено взаимодействие одиночных атомов с поверхностями~\cite{sipkens_examination_2015} и кластерами~\cite{insepov_kinetics_1991, insepov_jet_1991}. Эти исследования позволили определить такие ключевые параметры, как коэффициенты прилипания и испарения, и заложили основу для моделей кинетики конденсации в пересыщенном паре~\cite{insepov_kinetics_1991, insepov_jet_1991}.

Более сложные подходы включают статистическое моделирование систем, содержащих множество атомов (например, Cu и Ar), где процесс роста кластеров описывается в терминах классической теории нуклеации~\cite{eremin_influence_2008, vorontsov_statistics_2011, vorontsov_kinetics_2012, korenchenko_analysis_2016}. Такие работы позволяют изучать не только скорости роста, но и энергетические распределения в системе, включая эффективность отвода тепла, непосредственно связанную с КТА.

Коэффициент термической аккомодации является фундаментальным параметром, связывающим атомистическую картину взаимодействия с макроскопическим коэффициентом теплоотдачи~\cite{kurganov_thermal_2016, minakov_thermal_2015, popov_thermal_2016, valueva_nummodelling_2014}. Влияние наночастиц на эффективную теплопроводность сред также активно исследуется, в том числе и методами вычислительной гидродинамики~\cite{timofeeva_particle_2010}.

Альтернативой ресурсоемкому МД-моделированию являются кинетические подходы, основанные на решении уравнений для функции распределения кластеров по размерам, которые позволяют изучать эволюцию системы в условиях сильного и слабого пересыщения~\cite{goncharov_modelling_2011}. Кроме того, продолжаются работы по модификации и верификации классической теории нуклеации для объяснения экспериментальных данных, получаемых в различных условиях~\cite{fisenko_microstructure_2013}. Подробный обзор современных представлений о механизмах образования кластеров приведен в~\cite{varaksin_clusterization_2014}.





